\chapter{Literature Review and Existing Approaches}\label{ch:litreview}

This chapter summarises the published techniques the platform builds on. Statements about what this project does
are design choices of the team and are identified as such; they are justified by measurements in
\autoref{ch:testing}, not by the cited works.

\section{Paper Discovery and Metadata}

arXiv exposes an Atom-based query API that returns identifiers, versions, titles, authors, abstracts, categories and
links to PDF files \cite{arxivapi}. Identifiers are versioned (for example \code{1706.03762v7}), so a corpus can be
pinned to exact versions, which this project uses for its evaluation corpus. The API documentation asks clients to
space requests; the project observed HTTP 429 responses when this was not respected and added rate limiting and
retries accordingly (\autoref{ch:testing}).

\section{PDF Processing and Chunking}

Scientific PDFs encode text as positioned glyphs on pages. Libraries such as PyMuPDF extract text per page
\cite{pymupdf}, which makes it possible to record the page on which every character lies. Because embedding models
accept bounded inputs, extracted text is split into chunks. Fixed-size windows with overlap are a common baseline;
the overlap reduces the chance that a sentence is split across chunk boundaries. A known weakness of plain-text
extraction is that tables lose their row and column structure, which the project observed in practice
(\autoref{ch:results}).

\section{Embeddings and Semantic Search}

Transformer encoders \cite{vaswani2017attention,devlin2019bert} can be fine-tuned to map sentences to vectors
whose cosine similarity reflects semantic similarity. Sentence-BERT introduced siamese fine-tuning of BERT for this
purpose \cite{reimers2019sbert}; the \code{all-MiniLM-L6-v2} model used here is a compact sentence encoder from
that family \cite{minilmcard,sbertdocs}, based on the MiniLM distillation approach \cite{wang2020minilm}, and
produces 384-dimensional vectors. Dense passage retrieval showed that learned dense representations can outperform
lexical retrieval for open-domain question answering \cite{karpukhin2020dpr}.

\section{Vector Storage and Approximate Nearest-Neighbour Search}

Exact nearest-neighbour search scans every vector. Graph-based indexes such as Hierarchical Navigable Small World
(HNSW) graphs trade a small loss of recall for sub-linear query time \cite{malkov2020hnsw}. The pgvector extension
adds a vector type, distance operators and HNSW indexes to PostgreSQL \cite{pgvector}, which lets one database hold
papers, chunks, embeddings and question history with ordinary transactions and constraints. At the size of the
project's corpus the PostgreSQL planner chooses an exact scan rather than the HNSW index (\autoref{ch:testing}).

\section{Lexical Retrieval and Rank Fusion}

BM25 ranks documents by term frequency and inverse document frequency with length normalisation
\cite{robertson2009bm25,manning2008ir}. PostgreSQL provides full-text search with stemming and a cover-density
ranking function, \code{ts\_rank\_cd} \cite{pgtextsearch}; the project uses it as a BM25-style lexical ranker, not
as an exact BM25 implementation. Lexical and dense rankings fail on different queries, so combining them is
attractive. Reciprocal rank fusion (RRF) scores each document by $\sum_r 1/(k + \mathrm{rank}_r)$ over the input
rankings; its authors reported that it outperformed Condorcet fusion and individual learning-to-rank methods
\cite{cormack2009rrf}.

\section{Re-ranking}

A cross-encoder reads the query and a candidate passage together and outputs a relevance score. BERT-based
cross-encoders substantially improved passage ranking on MS MARCO \cite{nogueira2019rerank}. Because the query must
be re-encoded with every candidate, a cross-encoder is applied only to a short candidate list. The project uses the
publicly available \code{cross-encoder/ms-marco-MiniLM-L-6-v2} model \cite{crossencodercard}.

\section{Retrieval-Augmented Generation and Grounding}

RAG conditions a generator on retrieved passages \cite{lewis2020rag}; surveys describe many variants of retrieval,
context construction and evaluation \cite{gao2023ragsurvey}. Two findings from the literature shaped the design.
First, models do not use long contexts uniformly and can miss information in the middle of the context
\cite{liu2024lostmiddle}, which motivates bounded, ranked context. Second, generated text may contain claims not
supported by the input \cite{ji2023hallucination}. The project therefore requires the model to cite passage labels
for every claim, resolves those labels on the server, and treats a structurally valid citation as necessary but
not sufficient evidence of support.

\section{Evaluating RAG Systems}

Retrieval is commonly evaluated with Recall@$k$ and mean reciprocal rank (MRR) \cite{manning2008ir,voorhees1999trec8qa}.
Answer quality is harder to measure. Reference-free frameworks such as RAGAS score faithfulness and relevance with
a language model \cite{es2023ragas}, and studies of ``LLM-as-a-judge'' report useful agreement with humans together
with biases \cite{zheng2023judge}. The project uses a local model as a supplementary judge only and reports its
disagreements with manual inspection (\autoref{ch:testing}).

\section{Prompt Injection}

Text retrieved into a model's context can contain instructions. Indirect prompt injection showed that such text
can take control of LLM-integrated applications \cite{greshake2023injection}, and the OWASP Top 10 for LLM
applications lists prompt injection first \cite{owaspllm}. The project treats retrieved text as untrusted data: it is
delimited, delimiter-like text is neutralised, the instructions forbid following it, and the output format is
constrained by a JSON schema.

\section{Local Language Models}

Open-weight models can be served locally; Ollama provides a local HTTP API for running such models
\cite{ollama}. The project uses \code{qwen2.5:3b} from the Qwen2.5 family \cite{qwen2024qwen25}, chosen because it
fits in the 4\,GB of GPU memory on the reference laptop. Local inference keeps documents private but limits model
size, which affects answer quality and latency.

\section{Summary of Design Choices}

\autoref{tab:lit-choices} relates the main design choices of the team to the literature that motivates them.

\begin{table}[htbp]
  \centering
  \caption{Design choices and their background}\label{tab:lit-choices}
  \begin{tabularx}{\textwidth}{@{}Lp{4.6cm}@{}}
    \toprule
    \textbf{Choice in this project} & \textbf{Background} \\
    \midrule
    MiniLM sentence embeddings, cosine similarity & \cite{reimers2019sbert,wang2020minilm,karpukhin2020dpr} \\
    pgvector inside PostgreSQL & \cite{pgvector,malkov2020hnsw} \\
    Full-text ranking fused with dense ranking (RRF, $k=60$) & \cite{robertson2009bm25,cormack2009rrf,pgtextsearch} \\
    Cross-encoder re-ranking of a short candidate list & \cite{nogueira2019rerank} \\
    Answers only from retrieved, labelled passages; server-side citation checks & \cite{lewis2020rag,ji2023hallucination} \\
    Bounded context (six passages by default) & \cite{liu2024lostmiddle} \\
    Retrieved text treated as untrusted & \cite{greshake2023injection,owaspllm} \\
    Recall@$k$, MRR; local judge as a supplement only & \cite{manning2008ir,voorhees1999trec8qa,zheng2023judge} \\
    \bottomrule
  \end{tabularx}
\end{table}
